\chapter*{Final Considerations}
\addcontentsline{toc}{chapter}{Final Considerations}

The discipline of International Political Economy gained relevance in International Relations studies by, paradoxically, offering an approach to the International System that includes non-state actors. The British scholar Susan Strange inaugurated a vision of actors, activities, and individuals in the International System from a perspective that seeks to answer political-economic questions through four primary structures. This new approach includes non-state actors and brings academic studies closer to practical reality, since, although States are those ultimately responsible for sustaining order and security within national territories, they cannot be considered the only authorities that affect people and services.

Although not a theorist, Strange (1988) was a scholar and erudite who provided a vision of the operations of the international political-economic fabric in a way that encompasses four structures that, according to her, lie at the core of any society, whether at the global or local level: security, production, finance, and knowledge. After all, all minimally organized societies concern themselves with security and survival; with the production of food, tools, and products in general; with the creation of credit; and with the accumulation of knowledge and information that can be passed on between generations. The international system is observed, therefore, from an analytical framework that allows the participation of actors beyond the State -- since these are not the only ones that affect the four structures and people. Strange's IPE paradigm aims to understand the way in which power operates in the socio-political-economic fabric and, for this reason, she directed criticism at both economists (for ignoring the element of power in their analyses) and political scientists (for ignoring important aspects of the economy). This new paradigm seeks to provide sufficient resources for analyses that cover both segments within a global, or rather, international context.

It is from the IPE paradigm addressed in the 1988 work that Strange gathers the foundations that underpin her vision of the diffusion of power, affirming that State power has declined, in a qualitative sense, due to the rise of other authorities acting globally. According to her, such authorities have sufficient power not only to allocate basic values of human organization in society, but also sufficient capacity to alter outcomes and redefine options for others. These diverse non-state authorities are broadly embedded within the four primary structures. Their indirect power, or rather, structural power, is precisely the result of the action of structural dynamics, while direct power is well known in the discipline of International Relations: also known as relational power.

In this research, we sought to analyze the diffusion of power through the low-latency anonymous network ``The Onion Router.'' This network operates through the Internet, being part of the set of networks that form the ``Dark Web'' -- a region of the web where there is an active effort to maintain the anonymity of communications and users. This anonymous network seeks to be a relevant communication channel in the face of digital surveillance perpetrated by governmental and private entities. As a means of communication, it is situated within the knowledge structure, one of the primary structures initially defined by Strange in 1988. According to her, in addition to encompassing beliefs, principles, and moral conclusions, this structure also encompasses what is known and understood, and extends over the communication channels through which information is communicated. For this reason, we can affirm that the low-latency anonymous network Tor is part of the knowledge structure. And, as belonging to the knowledge structure, one verifies the existence of subjects, also positioned within this structure, who operate on the anonymous network. If Strange (1988) indicates that the structural analytical framework can be understood from both a global and local perspective, there is no reason to suppose that it would be unfeasible to use this framework for observations in the cyber domain. After all, different state and non-state actors operate in the cyber environment, affecting people and services offered in this domain, even though she herself did not address this domain specifically.

This dissertation sought to answer the research question in order to verify in what way, and to what extent, the anonymous network The Onion Router contributes to the diffusion of power in the knowledge structure. As seen, both the knowledge structure and the phenomenon of the diffusion of power were elements defined by Strange (1988, 1996) and, for this reason, our research was based on her works during the investigation of the phenomenon in this specific communication channel of the cyber domain. After considerable examination of the works, we identified two relevant points that bear on our analysis: first, the finding that the diffusion of power operates in three dimensions; second, even though she did not address the cyber domain, Strange indicated the importance of the technological revolutions that were occurring at the end of the 1980s and 1990s.

The first dimension is that of ``authority.'' In this dimension, we sought to demonstrate that Strange's considerations regarding power are treated, at times, as considerations regarding authority. However, although the differences between ``power'' and ``authority'' are subtle, we believe there are sufficient indicators in the author's works to affirm that these two elements are not synonymous. Authority can be viewed in two ways: as an entity that gathers the trust of peers and stakeholders upon itself; and also as the very exercise of power. In the first instance, we refer to ``being an authority'' -- which can be achieved through transfer, grant, concession, delegation, etc. This authority, therefore, can arise formally or informally. In the second instance, we refer to ``exercising authority.'' This exercise of authority implies the use of power, since exercising authority means that the authority uses all means available to it to put desired ends into practice. In other words, authority exercises power in order to realize its will. We emphasize, however, that both forms of authority can only be tested, verified, and examined through the third dimension -- outcomes.

The second dimension refers to ``control.'' Initially, we verified that on different occasions in the analyzed works, Strange (1988) indicates the control of an object by an actor to denote that actor's power. We gathered the citations in which the author makes these assertions and identified both the actor in question and the object controlled by them, giving rise to ``Appendix 3'' of this dissertation. This gave us sufficient input to examine the role of control in conferring authority and power upon actors within a structure. In the case of the knowledge structure, the role of control proved to be essential since it lies at the core of the definition of power in the knowledge structure: power is attributed to those who occupy a prominent position, or a decision-making position, in the knowledge structure and falls upon (1) those recognized by society as holders of knowledge; (2) those responsible for its storage; (3) and those who control the channels through which information and knowledge are transmitted. For the purposes of analysis, in this dimension, we assumed that control can be none, partial, or absolute.

The third dimension refers to ``outcomes.'' In this dimension, outcomes reflect the alteration of the status quo or its maintenance. The outcome attests to, or refutes, the authority of an actor and/or their control over an object they consider pertinent to achieving the desired result. That is, we can assume that an actor has authority -- after all, ``being an authority'' demonstrates the trust of peers and stakeholders. This authority is not synonymous with power. However, at some point, the actor may exercise the authority invested in them (formally or informally) and alter the status quo in favor of a scenario desired by them. In this situation, one verifies that the authority, in fact, has power. But if the actor, upon exercising authority (or rather, power), does not bring about the alteration of the status quo in favor of a scenario desired by them, we cannot admit that the actor still holds authority. Authority, in fact, can only be determined on the basis of this dimension, of outcomes. Similarly, the relevance of an object for a given actor in the pursuit of a specific result can only be determined if control of that object is pertinent to the alteration of the status quo. If it is not, the actor's control of that object does not result in power.

Our investigation was limited to examining the diffusion of power in the domain of the anonymous network ``The Onion Router.'' This investigation took place through the elaboration of a database containing occurrences derived from journalistic articles. These articles originated from the ten largest newspapers by number of visitors per click, as presented by Comscore (2012). The means by which the articles were selected to compose the database was presented in the third chapter of this dissertation. From this composition, we were able to make some interpretations regarding, for example, the popularity of the Tor network from 2013 onward. Furthermore, we carried out analyses of the diffusion of power in three dimensions according to three distinct segments: first, we presented a table with the exposition of the 25 actors (originating from the journalistic articles); second, the actors whose news content was similar were grouped into the same ``thematic group'' and we made considerations about the diffusion for the groups; and third, we examined the diffusion of power through the total occurrences between the years 2007--2017. Through our analysis, we were able to verify that the diffusion of power occurs for the majority of verified occurrences. There are sufficient indications to interpret that the specific technological nature of the Tor anonymous network is a preponderant factor for the presence of the phenomenon in the knowledge structure in the 21st century.

Lastly, we also brought important considerations regarding the phenomenon of the diffusion of power under Nye's paradigm in relation to ``cyberpower'' that originates from cyberspace. Nye's discourse on power resembles Strange's discourse in that it concerns a ``vertical'' phenomenon, in which it is supposed that power is ``diluting'' in the global political and economic fabric from the State toward non-state actors. Nye also specifically affirms the existence of a cyber power (cyberpower), becoming the first major theorist in the field of International Relations to make considerations about power in the new domain of the 21st century. It is possible to perceive that the definition of cyber power, coined by Nye (2011), is intimately related to a set of resources that derive from the knowledge structure -- albeit indirectly.

Finally, we presented charts, tables, and figures whose purpose was to serve as support for our analyses and to present the ``findings'' made during this present research.
